\documentclass[12pt]{article}
\usepackage{graphicx}

\usepackage{hyperref}  %allows for clickable reference
\hypersetup{colorlinks = true,citecolor = blue, linkcolor = blue, urlcolor = blue}
\usepackage{subcaption}

\usepackage{amsmath}
\usepackage{pifont}
\usepackage{wrapfig}
\usepackage{multicol}
\setlength{\columnsep}{1.5cm}
\setlength{\columnseprule}{1pt}
 
\usepackage{tikz}
\usetikzlibrary{calc}


\usepackage{fancyhdr}

\lhead{ }
\rhead{ }
\chead{Assignment-2}
 
\renewcommand{\headrulewidth}{0.4pt} %thickness of Line
\rfoot{Page \thepage}
\fancyfoot[CO,RE]{Vivek Kumar}
\fancyfoot[LO,CE]{10-Feb-2022}
\renewcommand{\footrulewidth}{0.2pt}  %thickness of line

 
\begin{document}
\pagestyle{empty}

\begin{tikzpicture}
[remember picture, overlay]
x\draw [line width=2pt]($(current page.north west)+(1in,-1in)$) rectangle ($(current page.south east)+(-1in, 1in)$);x

\end{tikzpicture}


   \begin{center}
       \vspace{1cm}
	   \Large
       \textbf{National Institute Of Technology, Raipur } 
       \vspace{1.5cm}
     
       \includegraphics[scale=0.1]{nit-raipur-logo.png}\\
       \vspace{0.8cm}
       \Huge
       Assignment-3\\
       \vspace{0.8cm}
      \textbf{ "Two page write-up on Future of Healthcare."}
      
	\vfill      
      
   \begin{multicols}{2} 
   \begin{flushleft}
       \large
       \textbf{Submitted by:}\\
       Name: Vivek Kumar\\
       Roll no: 21111071\\
       Branch : Biomedical\\
       Semester : $1^{st}$ \\
       NIT Raipur,Chhattisgarh\\
       \columnbreak
       \textbf{Under the Supervision of:}\\
       Mr.Saurabh Gupta
       Department of Biomedical Engineering
       NIT Raipur, Chhattisgarh
    \end{flushleft}
    \end{multicols}    
            
   \end{center}

\clearpage
\pagestyle{fancy}
\tableofcontents
\clearpage
\section{Overview}
Health Care is One of the India's Largest Sector both in terms of Revenue and Employement. Healthcare is a constantly evolving Sector. This evolution will takes many different forms and moves well beyond new medical discoveries. Investments are being made on forward-thinking innovations.
\textbf{Indian Healthcare} is expected to reach \$372 Billion by 2022 growing with a CAGR of over 16\%.
\section{Discussion}
The Covid-19 crisis showed us that even in an industry like healthcare, the change could be brought and accepted when needed. Many people avoided clinics and hospitals in fear of catching the virus, resulting in quicker opportunities for creativity and innovation.\\

\subsection{Analysis of Current Healthcare Industry}

According to \textbf{\href{https://www.himss.org/resources/future-healthcare-report-presentation}{HIMSS 2022 Report on Future of Healthcare}} on Hospitals, Health Organisations and Health Systems:-
\begin{itemize}

\item[\ding{234}] 80\% Plan to increase their digital investments in the next Five Years.

\item[\ding{234}] Organizations adopting AI/ML will see an uptick in clinical use cases compared to today's largely operational focus of AI/ML, 76\%believe that investments will increase in AI and digital health in the next five years. They see value in the business case to invest in AI and digital health.

\item[\ding{234}] There's a focus on expansion of value-based care (VBC) innovation and adoption to keep up with increased provider adoption and government policies.

\item[\ding{234}] Personalized care offerings are expected to increase.

\item[\ding{234}] Recent developments such as telehealth and EHRs(Electronic Health Record) allow patients to have easier access to care and increased information regarding the state of their health. Innovations such as wearable technology help patients to boost their health and well-being.

\end{itemize}


\subsection{Key Future Technologies}
\vspace{5mm}

\begin{itemize}
\large
 \item[\ding{110}]\textbf{\underline{In The Next 5-10 Years}}
\normalsize
 \item[\ding{234}] Portable Medical Devices for Professionals.
 
 \item[\ding{234}] Health Sensors for Consumers : Wearable Devices, Digital Tattoos, Smart Clothes for Health Monitoring.
 
 \item[\ding{234}] IoT in Health Care : Internet of Things HealthCare Devices.
 
 \item[\ding{234}] 3-D Printing : 3-D Printed Tissues, Bones, Organs.
 
 \item[\ding{234}] Robotics : Enhance the mobility of the Patients and may Provide Emotional Support.
 
 \item[\ding{234}] Virtual Reality in Health Care : Will be used Increasingly in Medical Training and Speeding up Recovery of Patients.
 
 \item[\ding{234}] Genomic Analysis : Receiving Treatment Customised to Our Molecular Make-Up and Genome Background.
 
 \item[\ding{234}] AI and ML in Medical Decision Support.
 
\vspace{5mm}
\large
\item[\ding{110}]\textbf{\underline{Beyond 5-10 years}}
\normalsize
\item Precision Medicine : to predict more accurately which treatment and prevention strategies for a particular disease will work Using Big Data and AI.

\item Augmented Reality : Will be Used for Diagnosis and Health Education.

\item Tele-Medicine : Will Solve  the Lack of Doctors Shortage
\item Devices inside the body (Nano-Technology): Will provide Real Time 24hr Monitoring.

\item Man made Organs : 3-D Printing of Organs.

\item Ageing Research : Ageing will become a Treatable Disease and People will cease to age.

\item Designer Babies : Manipulate DNAs to Customise Babies.

\item BCI(Brain-Computer Interface) Technology : Our Brain will be Directly Connected to AI.

\end{itemize}
\subsection{Hurdles Behind Healthcare}
\begin{itemize}
 
\item[\ding{111}] Rising Cost and varying Quality : It is One of the Main Challange in the way of Affordable Healthcare System. The treatment and Services might be more expensive in future.

\item[\ding{111}] Infectious diseases : Sometimes We find us so helpless to recognize a new disease and to design the treatment. We need time to understand the nature of the virus. During this period, it will cost the lives of many.

\item[\ding{111}] Data Privacy : Data privacy should be top of mind and not an afterthought. After all, data privacy protects the patient by protecting his or her data.

\item[\ding{111}] Healthcare supply chain management : Supply chain deficiencies made it difficult for some providers to deliver the best care possible.
 
\end{itemize}
 
\section{Conclusion}
It is clear that in addition to more efficient billing, less wasteful tests and procedures will be done as better information on appropriate care becomes available, and more efficient care models will emerge with technology for care of patients at home, these improvements will be dwarfed by increasing costs. In the next 10 years, process and outcomes will be optimized for a significant proportion of patients with relatively common disease smaking it possible to develop a better understanding of the best care delivery model and this will improve the patient experience.  The ultimate goal is to improve patient value.
\end{document}